% ============================================================
% content/questions.tex —— 2025 湖南卷物理 · 试题
% 来源：2025/2025_物理真题_含解析.md（官方 Word 版转录）
% 题号保留官方编号；图形见 figures/（TikZ 重绘）与 photos/（非矢量例外）
% ============================================================

\gkpart{一、选择题}

\textbf{（一）单项选择题（本题共 6 小题，每小题 4 分，共 24 分。在每小题给出的四个选项中，只有一项是符合题目要求的。）}

\begin{question}
\gkti{1}关于原子核衰变，下列说法正确的是\gkblank

\begin{gkopts}(1)
  \task A．原子核衰变后生成新核并释放能量，新核总质量等于原核质量
  \task B．大量某放射性元素的原子核有半数发生衰变所需时间，为该元素的半衰期
  \task C．放射性元素的半衰期随环境温度升高而变长
  \task D．采用化学方法可以有效改变放射性元素的半衰期
\end{gkopts}
\end{question}

\begin{question}
\gkti{2}如图所示，物块以某一初速度滑上足够长的固定光滑斜面，物块的水平位移、竖直位移、水平速度、竖直速度分别用~$x$、$y$、$v_{x}$、$v_{y}$~表示。物块向上运动过程中，下列图像可能正确的是\gkblank

\gkfig{figures/q2-graphs}
\end{question}

\begin{question}
\gkti{3}如图所示，$ABC$~为半圆柱体透明介质的横截面，$AC$~为直径，$B$~为$\overset{\frown}{ABC}$的中点。真空中一束单色光从~$AC$~边射入介质，入射点为~$A$~点，折射光直接由~$B$~点出射。不考虑光的多次反射，下列说法正确的是\gkblank

\gkfig{figures/q3-semicircle}

\begin{gkopts}(2)
  \task A．入射角 $\theta$ 小于 $45^{\circ}$
  \task B．该介质折射率大于 $\sqrt{2}$
  \task C．增大入射角，该单色光在 $\overset{\frown}{BC}$ 上可能发生全反射
  \task D．减小入射角，该单色光在 $\overset{\frown}{AB}$ 上可能发生全反射
\end{gkopts}
\end{question}

\begin{question}
\gkti{4}我国研制的“天问二号”探测器，任务是对伴地小行星及彗星交会等进行多目标探测。某同学提出探究方案，通过释放卫星绕小行星进行圆周运动，可测得小行星半径~$R$~和质量~$M$。为探测某自转周期为~$T_{0}$~的小行星，卫星先在其同步轨道上运行，测得距离小行星表面高度为~$h$，接下来变轨到小行星表面附近绕其做匀速圆周运动，测得周期为~$T_{1}$。已知引力常量为~$G$，不考虑其他天体对卫星的引力，可根据以上物理量得到
\[
  R=\frac{a^{\frac{2}{3}}}{b^{\frac{2}{3}}-a^{\frac{2}{3}}}\,h,\qquad
  M=\frac{4\pi^{2}R^{3}}{Gc^{2}}.
\]
下列选项正确的是\gkblank

\gkfig{figures/q4-orbit}

\begin{gkopts}(1)
  \task A．$a$ 为 $T_{1}$，$b$ 为 $T_{0}$，$c$ 为 $T_{1}$
  \task B．$a$ 为 $T_{1}$，$b$ 为 $T_{0}$，$c$ 为 $T_{0}$
  \task C．$a$ 为 $T_{0}$，$b$ 为 $T_{1}$，$c$ 为 $T_{1}$
  \task D．$a$ 为 $T_{0}$，$b$ 为 $T_{1}$，$c$ 为 $T_{0}$
\end{gkopts}
\end{question}

\begin{question}
\gkti{5}如图所示，两带电小球的质量均为~$m$，小球~$A$~用一端固定在墙上的绝缘轻绳连接，小球~$B$~用固定的绝缘轻杆连接。小球~$A$~静止时，轻绳与竖直方向的夹角为 $60^{\circ}$，两球连线与轻绳的夹角为 $30^{\circ}$，整个系统在同一竖直平面内，重力加速度大小为~$g$。下列说法正确的是\gkblank

\gkfig{figures/q5-balls}

\begin{gkopts}(2)
  \task A．小球 $A$ 静止时，轻绳上拉力为 $2mg$
  \task B．小球 $A$ 静止时，小球 $A$ 与小球 $B$ 间的库仑力为 $2mg$
  \task C．若将轻绳剪断，则剪断瞬间小球 $A$ 加速度大小为 $g$
  \task D．若将轻绳剪断，则剪断瞬间轻杆对小球 $B$ 的作用力变小
\end{gkopts}
\end{question}

\begin{question}
\gkti{6}如图所示，某小组设计了灯泡亮度可调的电路，$a$、$b$、$c$~为固定的三个触点，理想变压器原、副线圈匝数比为~$k$，灯泡~$L$~和三个电阻的阻值均恒为~$R$，交流电源输出电压的有效值恒为~$U$。开关~$S$~与不同触点相连，下列说法正确的是\gkblank

\gkfig{figures/q6-circuit}

\begin{gkopts}(1)
  \task A．$S$ 与 $a$ 相连，灯泡的电功率最大
  \task B．$S$ 与 $a$ 相连，灯泡两端的电压为 $\dfrac{kU}{k^{2}+3}$
  \task C．$S$ 与 $b$ 相连，流过灯泡的电流为 $\dfrac{U}{(k^{2}+2)R}$
  \task D．$S$ 与 $c$ 相连，灯泡的电功率为 $\dfrac{U^{2}}{(k^{2}+1)R}$
\end{gkopts}
\end{question}

\textbf{（二）多项选择题（本题共 4 小题，每小题 5 分，共 20 分。在每小题给出的四个选项中，有多项符合题目要求；全部选对的得 5 分，选对但不全的得 3 分，有选错的得 0 分。）}

\begin{question}
\gkti{7}如图所示，$A(0,\,0)$、$B(4,\,0)$、$C(0,\,3)$~在~$xOy$~平面内，两波源分别置于~$A$、$B$~两点。$t=0$~时，两波源从平衡位置起振，起振方向相同且垂直于~$xOy$~平面，频率均为~\SI{2.5}{\hertz}。两波源持续产生振幅相同的简谐横波，波分别沿~$AC$、$BC$~方向传播，波速均为~\SI{10}{\metre\per\second}。下列说法正确的是\gkblank

\gkfig{figures/q7-waves}

\begin{gkopts}(2)
  \task A．两横波的波长均为 \SI{4}{\metre}
  \task B．$t=\SI{0.4}{\second}$ 时，$C$ 处质点加速度为 0
  \task C．$t=\SI{0.4}{\second}$ 时，$C$ 处质点速度不为 0
  \task D．$t=\SI{0.6}{\second}$ 时，$C$ 处质点速度为 0
\end{gkopts}
\end{question}

\begin{question}
\gkti{8}一匀强电场的方向平行于~$xOy$~平面，平面内~$A$~点和~$B$~点的位置如图所示。电荷量为 $+q$、$-q$~和 $+2q$~的三个试探电荷先后分别置于~$O$~点、$A$~点和~$B$~点时，电势能均为~$E_{p}$（$E_{p}>0$）。下列说法正确的是\gkblank

\gkfig{figures/q8-field}

\begin{gkopts}(1)
  \task A．$OA$ 中点的电势为零
  \task B．电场的方向与 $x$ 轴正方向成 $60^{\circ}$ 角
  \task C．电场强度的大小为 $\dfrac{\sqrt{2}\,E_{p}}{qd}$
  \task D．电场强度的大小为 $\dfrac{2\sqrt{2}\,E_{p}}{qd}$
\end{gkopts}
\end{question}

\begin{question}
\gkti{9}如图所示，关于~$x$~轴对称的光滑导轨固定在水平面内，导轨形状为抛物线，顶点位于~$O$~点。一足够长的金属杆初始位置与~$y$~轴重合，金属杆的质量为~$m$，单位长度的电阻为~$r_{0}$。整个空间存在竖直向上的匀强磁场，磁感应强度为~$B$。现给金属杆一沿~$x$~轴正方向的初速度~$v_{0}$，金属杆运动过程中始终与~$y$~轴平行，且与电阻不计的导轨接触良好。下列说法正确的是\gkblank

\gkfig{figures/q9-parabola}

\begin{gkopts}(1)
  \task A．金属杆沿 $x$ 轴正方向运动过程中，金属杆中电流沿 $y$ 轴负方向
  \task B．金属杆可以在沿 $x$ 轴正方向的恒力作用下做匀速直线运动
  \task C．金属杆停止运动时，与导轨围成的面积为 $\dfrac{mv_{0}r_{0}}{B^{2}}$
  \task D．若金属杆的初速度减半，则金属杆停止运动时经过的距离小于原来的一半
\end{gkopts}
\end{question}

\begin{question}
\gkti{10}如图所示，某爆炸能量测量装置由装载台和滑轨等构成，$C$~是可以在滑轨上运动的标准测量件，其规格可以根据测量需求进行调整。滑轨安装在高度为~$h$~的水平面上。测量时，将弹药放入装载台圆筒内，两端用物块~$A$~和~$B$~封装，装载台与滑轨等高。引爆后，假设弹药释放的能量完全转化为~$A$~和~$B$~的动能。极短时间内~$B$~嵌入~$C$~中形成组合体~$D$，$D$~与滑轨间的动摩擦因数为~$\mu$。$D$~在滑轨上运动~$s_{1}$~距离后抛出，落地点距抛出点水平距离为~$s_{2}$，根据~$s_{2}$~可计算出弹药释放的能量。某次测量中，$A$、$B$、$C$~质量分别为 $3m$、$m$、$5m$，$s_{1}=\dfrac{h}{\mu}$，整个过程发生在同一竖直平面内，不计空气阻力，重力加速度大小为~$g$。则\gkblank

\gkphoto[0.7\textwidth]{photos/image9-apparatus.png}{爆炸能量测量装置示意图（原图保留，非矢量例外）}

\begin{gkopts}(1)
  \task A．$D$ 的初动能与爆炸后瞬间 $A$ 的动能相等
  \task B．$D$ 的初动能与其落地时的动能相等
  \task C．弹药释放的能量为 $36mgh\left(1+\dfrac{s_{2}^{2}}{4h^{2}}\right)$
  \task D．弹药释放的能量为 $48mgh\left(1+\dfrac{s_{2}^{2}}{4h^{2}}\right)$
\end{gkopts}
\end{question}

\gkpart{二、实验题}

\textbf{（本题共 2 小题，共 16 分。）}

\begin{question}
\gkti{11}（7 分）某同学通过观察小球在黏性液体中的运动，探究其动力学规律，步骤如下：

\gkphoto[0.45\textwidth]{photos/image10-micrometer.png}{图 1（螺旋测微器，原图保留，非矢量例外）}

\gkphoto[0.4\textwidth]{photos/image11-strobe.png}{图 2（频闪摄影装置及小球下落照片示意图，原图保留，非矢量例外）}

\begin{enumerate}[leftmargin=2em, label=(\arabic*)]
  \item 用螺旋测微器测量小球直径 $D$ 如图 1 所示，$D=$\gkblank mm。
  \item 在液面处由静止释放小球，同时使用频闪摄影仪记录小球下落过程中不同时刻的位置，频闪摄影仪每隔 \SI{0.5}{\second} 闪光一次。装置及所拍照片示意图如图 2 所示（图中的数字是小球到液面的测量距离，单位是 cm）。
  \item 根据照片分析，小球在 $A$、$E$ 两点间近似做匀速运动，速度大小 $v=$\gkblank m/s（保留 2 位有效数字）。
  \item 小球在液体中运动时受到液体的黏滞阻力 $f=kDv$（$k$ 为与液体有关的常量），已知小球密度为 $\rho$，液体密度为 $\rho_{0}$，重力加速度大小为 $g$，则 $k$ 的表达式为 $k=$\gkblank（用题中给出的物理量表示）。
  \item 为了进一步探究动力学规律，换成直径更小的同种材质小球，进行上述实验，匀速运动时的速度将\gkblank（选填“增大”“减小”或“不变”）。
\end{enumerate}
\end{question}

\begin{question}
\gkti{12}（9 分）车辆运输中若存在超载现象，将带来安全隐患。由普通水泥和导电材料混合制成的导电水泥，可以用于监测道路超载问题。某小组对此进行探究。

\gkphoto[0.6\textwidth]{photos/image12-multimeter.png}{图 1（多用电表指针读数，原图保留，非矢量例外）}

\gkphoto[0.6\textwidth]{photos/image13-cement.png}{图 2（导电水泥块及部分接线，原图保留，非矢量例外）}

\begin{enumerate}[leftmargin=2em, label=(\arabic*)]
  \item 选择一块均匀的长方体导电水泥块样品，用多用电表粗测其电阻。将多用电表开关旋转到“$\times$1 k”挡，正确操作后，指针位置如图 1 所示，则读数为\gkblank\ \si{\ohm}。
  \item 为进一步提高实验精度，使用伏安法测量水泥块电阻，电源 $E$ 电动势 \SI{6}{\volt}，内阻可忽略，电压表量程 $0\sim$\SI{6}{\volt}，内阻约 \SI{10}{\kilo\ohm}，电流表量程 $0\sim$\SI{600}{\micro\ampere}，内阻约 \SI{100}{\ohm}。实验中要求滑动变阻器采用分压接法，在图 2 中完成余下导线的连接。
  \item 如图 2 所示，测量水泥块的长为 $a$，宽为 $b$，高为 $c$。用伏安法测得水泥块电阻为 $R$，则电阻率 $\rho=$\gkblank（用 $R$、$a$、$b$、$c$ 表示）。
  \item 测得不同压力 $F$ 下水泥块的电阻 $R$，算出对应的电阻率 $\rho$，作出 $\rho$--$F$ 图像如图 3 所示。
  \item 基于以上结论，设计压力报警系统，电路如图 4 所示。报警器在两端电压大于或等于 \SI{3}{\volt} 时启动，$R_{1}$ 为水泥块，$R_{2}$ 为滑动变阻器，当 $R_{2}$ 的滑片处于某位置，$R_{1}$ 上压力大于或等于 $F_{0}$ 时，报警器启动。报警器应并联在\gkblank（选填“$R_{1}$”或“$R_{2}$”）两端。
  \item 若电源 $E$ 使用时间过长，电动势变小，$R_{1}$ 上压力大于或等于 $F_{1}$ 时，报警器启动，则 $F_{1}$\gkblank（选填“大于”“小于”或“等于”）$F_{0}$。
\end{enumerate}

\gkfig{figures/q12-rhof}

\gkfig{figures/q12-alarm}
\end{question}

\gkpart{三、解答题}

\textbf{（本题共 3 小题，共 40 分。解答应写出必要的文字说明、方程式和重要的演算步骤，只写出最后答案的不能得分。）}

\begin{question}
\gkti{13}（10 分）用热力学方法可测量重力加速度。如图所示，粗细均匀的细管开口向上竖直放置，管内用液柱封闭了一段长度为~$L_{1}$~的空气柱。液柱长为~$h$，密度为~$\rho$。缓慢旋转细管至水平，封闭空气柱长度为~$L_{2}$，大气压强为~$p_{0}$。

\gkfig{figures/q13-tube}

\begin{enumerate}[leftmargin=2em, label=(\arabic*)]
  \item 若整个过程中温度不变，求重力加速度 $g$ 的大小。
  \item 考虑到实验测量中存在各类误差，需要在不同实验参数下进行多次测量，如不同的液柱长度、空气柱长度、温度等。某次实验测量数据如下：液柱长 $h=\SI{0.2000}{\metre}$，细管开口向上竖直放置时空气柱温度 $T_{1}=\SI{305.7}{\kelvin}$，水平放置时调控空气柱温度，当空气柱温度 $T_{2}=\SI{300.0}{\kelvin}$ 时，空气柱长度与竖直放置时相同。已知 $\rho=\SI{1.0e3}{\kilogram\per\cubic\metre}$，$p_{0}=\SI{1.0e5}{\pascal}$。根据该组实验数据，求重力加速度 $g$ 的值。
\end{enumerate}
\end{question}

\begin{question}
\gkti{14}（14 分）如图所示，直流电源的电动势为~$E_{0}$，内阻为~$r_{0}$，滑动变阻器~$R$~的最大阻值为 $2r_{0}$，平行板电容器两极板水平放置，板间距离为~$d$，板长为 $\sqrt{3}\,d$，平行板电容器的右侧存在方向垂直于纸面向里的匀强磁场。闭合开关~$S$，当滑片处于滑动变阻器中点时，质量为~$m$~的带正电粒子以初速度~$v_{0}$~水平向右从电容器左侧中点~$a$~进入电容器，恰好从电容器下极板右侧边缘~$b$~点进入磁场，随后又从电容器上极板右侧边缘~$c$~点进入电容器，忽略粒子重力和空气阻力。

\gkfig{figures/q14-capacitor}

\begin{enumerate}[leftmargin=2em, label=(\arabic*)]
  \item 求粒子所带电荷量 $q$。
  \item 求磁感应强度 $B$ 的大小。
  \item 若粒子离开 $b$ 点时，在平行板电容器的右侧再加一个方向水平向右的匀强电场，电场强度大小为 $\dfrac{4\sqrt{3}\,E_{0}}{3d}$，求粒子相对于电容器右侧的最远水平距离 $x_{m}$。
\end{enumerate}
\end{question}

\begin{question}
\gkti{15}（16 分）某地为发展旅游经济，因地制宜利用山体举办了机器人杂技表演。表演中，需要将质量为~$m$~的机器人抛至悬崖上的~$A$~点，图为山体截面与表演装置示意图。$a$、$b$~为同一水平面上两条光滑平行轨道，轨道中有质量为~$M$~的滑杆。滑杆用长度为~$L$~的轻绳与机器人相连。初始时刻，轻绳绷紧且与轨道平行，机器人从~$B$~点以初速度~$v$~竖直向下运动，$B$~点位于轨道平面上，且在~$A$~点正下方，$AB=1.2L$。滑杆始终与轨道垂直，机器人可视为质点且始终在同一竖直平面内运动，不计空气阻力，轻绳不可伸长，$\sin 37^{\circ}=0.6$，重力加速度大小为~$g$。

\gkphoto[0.55\textwidth]{photos/image18-robot.png}{山体截面与表演装置示意图（原图保留，非矢量例外）}

\begin{enumerate}[leftmargin=2em, label=(\arabic*)]
  \item 若滑杆固定，$v=\sqrt{gL}$，当机器人运动到滑杆正下方时，求轻绳拉力的大小；
  \item 若滑杆固定，当机器人运动到滑杆左上方且轻绳与水平方向夹角为 $37^{\circ}$ 时，机器人松开轻绳后被抛至 $A$ 点，求 $v$ 的大小；
  \item 若滑杆能沿轨道自由滑动，$M=km$，且 $k\ge 1$，当机器人运动到滑杆左上方且轻绳与水平方向夹角为 $37^{\circ}$ 时，机器人松开轻绳后被抛至 $A$ 点，求 $v$ 与 $k$ 的关系式及 $v$ 的最小值。
\end{enumerate}
\end{question}
